\documentclass[11pt]{article}
\usepackage[margin=1in]{geometry}
\usepackage[T1]{fontenc}
\usepackage{lmodern}
\usepackage{amsmath,booktabs,hyperref,microtype}
\hypersetup{hidelinks}
\makeatletter\let\tableinput\@@input\makeatother
% Generated; do not edit. Held-out validation vs CIPHER.
\newcommand{\ValAssumedEhrResidual}{+8.5}
\newcommand{\ValAssumedGamma}{0.60}
\newcommand{\ValAssumedP}{0.281}
\newcommand{\ValAssumedPharmaResidual}{-6.3}
\newcommand{\ValAssumedTV}{9.2}
\newcommand{\ValAttackVersion}{17.1}
\newcommand{\ValHighTotal}{128}
\newcommand{\ValModelledRecords}{88}
\newcommand{\ValObsEhr}{21.6}
\newcommand{\ValObsImg}{20.5}
\newcommand{\ValObsLab}{27.3}
\newcommand{\ValObsPharma}{30.7}
\newcommand{\ValTau}{7}
\newcommand{\ValUniformGamma}{0.35}
\newcommand{\ValUniformP}{0.502}
\newcommand{\ValUniformTV}{8.0}


\title{Do attack-derived models predict where patients are harmed?\\
A first validation of an ATT\&CK-based clinical-impact model against coded harm}
\author{}
\date{2026-09-20 (working note, version 0.1.0)}

\begin{document}
\maketitle

\begin{abstract}
Threat models built on MITRE ATT\&CK reason about attacker behaviour on IT, not
about which patients are harmed. We ask, for the first time to our knowledge,
whether an ATT\&CK-based clinical-impact model's predicted distribution of harm
across clinical services matches a corpus of \emph{real coded patient harm} from
hospital ransomware incidents (CIPHER: \ValHighTotal{} high-severity records). To
avoid circularity, the predicted profiles are built \emph{without} CIPHER --- a
uniform ATT\&CK-null and the model's Neprash-anchored ``assumed'' degradation ---
and CIPHER is held out as ground truth. The real harm profile is
pharmacy-\ and laboratory-heaviest (pharmacy \ValObsPharma\%, laboratory
\ValObsLab\%, EHR \ValObsEhr\%, imaging \ValObsImg\%), whereas the model's assumed
profile is EHR-heaviest: it over-predicts EHR by \ValAssumedEhrResidual{} points
and under-predicts pharmacy by \ValAssumedPharmaResidual{} points, and needs a
larger reporting-bias to reconcile ($\gamma=\ValAssumedGamma$) than even a uniform
guess ($\gamma=\ValUniformGamma$). Small sample size means neither profile is
formally rejected, so the honest conclusion is not "the model is wrong" but "the
model's clinical-service resolution is not validated by real harm, and points the
wrong way." The contribution is the validation itself, and the CIPHER overlay that
supplies the clinical resolution ATT\&CK lacks --- not a new statistical method.
\end{abstract}

\section{Design}
CIPHER codes real patient harms by clinical specialty, technical domain, severity
and onset timing. We map its four technical domains with a modelled counterpart
(Health Records~$\to$~EHR, Laboratory Systems~$\to$~laboratory, ePrescribing~$\to$
pharmacy, Imaging~$\to$~imaging) and take the high-severity (impact
$\ge\ValTau$) per-service counts as the observed profile (\ValModelledRecords{}
records; the large unmodelled residual is reported, not folded in). The predicted
profiles use only model inputs independent of CIPHER: a \emph{uniform} profile (the
honest ATT\&CK null, since ATT\&CK has no clinical-service signal) and the model's
\emph{assumed} degradation, set from the Neprash aggregate. This held-out design
pre-empts training-on-test.

\section{Comparison}
We report the total-variation distance to the observed profile, an exact
Monte-Carlo multinomial goodness-of-fit $p$-value (valid at small counts), and the
\emph{reconciling $\gamma$}: the least bounded per-domain underreporting that would
place the predicted profile inside CIPHER's partial-identification share intervals.
A small $\gamma$ means the model is close; a large $\gamma$ means a real gap that
only heavy reporting bias could explain.

\section{Result}
Neither predicted profile is rejected at $\tau=\ValTau$ (Monte-Carlo
$p=\ValUniformP$ uniform, $p=\ValAssumedP$ assumed): with \ValModelledRecords{}
records across four services the test has little power, which we state plainly
rather than read as agreement. The informative signals are directional
(Table~\ref{tab:valres}): the assumed profile over-predicts EHR by
\ValAssumedEhrResidual{} points and under-predicts pharmacy by
\ValAssumedPharmaResidual{} points, and its reconciling $\gamma=\ValAssumedGamma$
exceeds the uniform null's $\ValUniformGamma$ (Table~\ref{tab:valgof}). In other
words, the model's clinical-service differentiation does not beat a uniform guess
against real harm, and it emphasises the wrong service. This is the first such
calibration check for an ATT\&CK-based clinical model, and it motivates the CIPHER
overlay --- an empirical per-service specialty, severity and onset-timing layer ---
as the clinical resolution ATT\&CK cannot supply on its own.

\begin{table}[t]\centering
\caption{Fit of each CIPHER-independent predicted profile to the observed harm
profile at $\tau=\ValTau$.}\label{tab:valgof}
\begin{tabular}{@{}lrrrr@{}}\toprule
Predicted profile & Total variation (\%) & MC $p$ & $\chi^2$ $p$ & Reconciling $\gamma$\\\midrule
\tableinput table_validation_gof_rows.tex
\bottomrule\end{tabular}
\end{table}

\begin{table}[t]\centering
\caption{Assumed profile vs.\ observed harm, per service (percent), with CIPHER
mean clinical-impact score.}\label{tab:valres}
\begin{tabular}{@{}lrrrr@{}}\toprule
Service & Predicted (\%) & Observed (\%) & Residual (pp) & CIPHER mean impact\\\midrule
\tableinput table_validation_residual_rows.tex
\bottomrule\end{tabular}
\end{table}

\section{Limitations}
CIPHER is a convenience sample of reported harms: selection and reporting bias, no
denominator, English-language and severe-harm skew. The comparison tests
consistency with \emph{this corpus}, not population truth; the reconciling $\gamma$
makes the reporting-bias assumption explicit rather than hidden. Counts are small,
so power is low and no strong rejection is claimed. Coding is single-label and the
ATT\&CK$\to$service mapping is coarse; the unmodelled residual is large. This is a
validation and an overlay resource, honestly scoped --- not a new statistical
method and not a population estimate.

\paragraph{Reproducibility.} Predicted profiles are CIPHER-independent; the
Monte-Carlo test uses a fixed seed; all numbers are generated from committed
tables under a provenance manifest over the pinned ATT\&CK bundle and CIPHER.

\end{document}
